\documentclass[../../../include/open-logic-section]{subfiles}

\begin{document}

\olfileid{his}{set}{infinitesimals}
\olsection{అనంతసూక్ష్మాలు మరియు అవకలనం}

పదిహేడో శతాబ్దం చివర్లో న్యూటన్, లైబ్నిజ్ ఒకరికొకరు స్వతంత్రంగా కలనశాస్త్రాన్ని కనుగొన్నారు. కలనశాస్త్రానికి ప్రత్యేకంగా ముఖ్యమైన అనువర్తనాల్లో \emph{అవకలనం} ఒకటి. స్థూలంగా చెప్పాలంటే, ఒక బిందువు వద్ద ప్రమేయం యొక్క ``మార్పు రేటు'', లేదా వాలు, అనే భావనను ఇవ్వడమే అవకలనం లక్ష్యం.

ఈ భావనను స్పష్టంగా చూపడానికి ఒక మార్గం ఉంది. కింద నలుపు రంగులో చూపిన $f(x) = \nicefrac{x^2}{4} + \nicefrac{1}{2}$ ప్రమేయాన్ని చూడండి:
\begin{center}
	\begin{tikzpicture}[scale=1]
	\draw[->, gray] (-1,0) -- (4.25,0) node[right] {$x$};
	\draw[->, gray] (0,-0.25) -- (0,5.25) node[above] {$f(x)$};
	
	\foreach \x/\xtext in {1/1, 2/2, 3/3, 4/4}
	\draw[shift={(\x,0)}, gray] (0pt,2pt) -- (0pt,-2pt) node[below] {$\xtext$};
	
	\foreach \y/\ytext in {1/1, 2/2, 3/3, 4/4, 5/5}
	\draw[shift={(0,\y)}, gray] (2pt,0pt) -- (-2pt,0pt) node[left] {$\ytext$};
	
	\draw[oldiagcolorC] (0.5, 9/16) -- (3.5, 57/16) -- (3.5, 9/16)--cycle;
	\draw[oldiagcolorD] (.5, 9/16) -- (2.5, 33/16) -- (2.5, 9/16) -- cycle;
	\draw[oldiagcolorE] (.5, 9/16) -- (1.5, 17/16) -- (1.5, 9/16) -- cycle;
	\draw[thick] (-1,.75) parabola bend (0,0.5) (4,4.5);
	\end{tikzpicture}
\end{center}
$c =
\nicefrac{1}{2}$ వద్ద ప్రమేయం వాలును తెలుసుకోవాలనుకుందాం. మొదట ఆ బిందువు వద్ద వాలుకు దగ్గరగా ఉండే కర్ణం గల త్రిభుజాన్ని గీయండి; పై చిత్రంలోని ఎరుపు త్రిభుజం అటువంటిదే కావచ్చు. మన త్రిభుజం పాదం పొడవు $\beta$ అయితే, దాని ఎత్తు $f(\nicefrac{1}{2}+\beta) - f(\nicefrac{1}{2})$; కనుక కర్ణం వాలు:
\[
\frac{f(\nicefrac{1}{2}+\beta) - f(\nicefrac{1}{2})}{\beta}.
\]
అందువల్ల పాదం పొడవు~$3$ ఉన్న \tetoken{ఎరుపు}{!!{colorC}} త్రిభుజం వాలు కచ్చితంగా~$1$. మరింత చిన్నదైన, పాదం పొడవు~$2$ ఉన్న \tetoken{నీలం}{!!{colorD}} త్రిభుజం కర్ణం మరింత మంచి ఉజ్జాయింపునిస్తుంది; దాని వాలు $\nicefrac{3}{4}$. ఇంకా చిన్నదైన, పాదం పొడవు~$1$ ఉన్న \tetoken{ఆకుపచ్చ}{!!{colorE}} త్రిభుజం ఇంకా మెరుగైన ఉజ్జాయింపునిస్తుంది; దాని వాలు $\nicefrac{1}{2}$.

త్రిభుజాలు చిన్నవవుతున్న కొద్దీ ఉజ్జాయింపులు మెరుగవుతాయి. అందువల్ల ఇలా అనవచ్చు: పాదం పొడవు \emph{అనంతసూక్ష్మం} అయిన త్రిభుజం కర్ణం, $c =
\nicefrac{1}{2}$ వద్ద వాలునే ఇస్తుంది. ఈ విధంగా $c$ బిందువు వద్ద $f$ ప్రమేయం యొక్క (మొదటి) వ్యుత్పన్నానికి ఒక సూత్రాన్ని పొందుతాం:
\[
{f'}(c) = \frac{f(c+\beta) - f(c)}{\beta} \text{ ఇక్కడ $\beta$ అనంతసూక్ష్మం.}
\]
స్థూలంగా చెప్పాలంటే న్యూటన్, లైబ్నిజ్ చెప్పింది ఇదే.

అయితే వారు అలా చెప్పినందున, \emph{అనంతసూక్ష్మం} అంటే ఏమిటని అడగాలి. ఇక్కడ తీవ్రమైన సందిగ్ధం ఉంది. $\beta = 0$ అయితే $0$తో భాగించాలి కనుక $f'$ నిర్వచించబడదు. కానీ $\beta >
0$ అయితే మనకు వాలుకే బదులు దాని \emph{ఉజ్జాయింపు} మాత్రమే వస్తుంది.

ఈ సందేహం కాలానికి అనుచితమైన తరువాతి ప్రశ్న కాదు. న్యూటన్ అనుచరులను విమర్శిస్తూ బర్క్లీ ఇలా అన్నారు:
\begin{quote}
సంకేతాలు ఏదో ఒకదానిని లేదా ఏదీ కానిదానిని సూచించవచ్చని నేను అంగీకరిస్తాను; అందువల్ల మొదటి సంకేతరూపం $c + \beta$లో $\beta$ పెరుగుదలను గానీ శూన్యాన్ని గానీ సూచించి ఉండవచ్చు. కానీ దేనిని సూచిస్తుందని తీసుకున్నా, అదే అర్థానికి అనుగుణంగా స్థిరంగా వాదించాలి; రెండు అర్థాలను కలిపి కొనసాగకూడదు. అలా చేయడం స్పష్టమైన కుతర్కం. (\citealt[\S{}XIII]{Berkeley1734}; ముందు పాఠ్యానికి సరిపడేలా చరాలను మార్చాం)
\end{quote}
బర్క్లీ విమర్శ నుంచి అనంతసూక్ష్మ కలనశాస్త్రాన్ని సమర్థించడానికి, ``అనంతసూక్ష్మాలు'' అనే మాట కేవలం అలంకారప్రాయమని ఎవరైనా అనవచ్చు. \emph{నిజంగా చాలా చిన్న} త్రిభుజాన్ని తీసుకుంటే, స్పర్శరేఖకు \emph{సరిపడా మంచి} ఉజ్జాయింపు దొరుకుతుందని చెప్పవచ్చు. దీనికీ బర్క్లీ సమాధానం ఇచ్చారు: ఇంజినీరింగ్‌కు అది సరిపోవచ్చు; కానీ గణితశాస్త్రపు \emph{స్థాయిని} అది దెబ్బతీస్తుంది, ఎందుకంటే
\begin{quote}
మనకు \emph{in rebus mathematicis errores qu\`{a}m minimi
non sunt contemnendi} అని చెప్పారు. [గణిత విషయాలలో అతి చిన్న తప్పులనూ నిర్లక్ష్యం చేయరాదు.] \citep[\S{}IX]{Berkeley1734}
\end{quote}
ఇటాలిక్‌లోని మాటలు న్యూటన్ స్వంత \emph{Quadrature of Curves} (1704) నుంచి దాదాపు యథాతథంగా తీసుకున్నవి.

బర్క్లీ తాత్విక అభ్యంతరాలు ఎంతో లోతైనవి. అయినప్పటికీ కలనశాస్త్రం అత్యంత విజయవంతమైంది; గణితశాస్త్రజ్ఞులు తప్పుల్లో పడకుండా దాన్ని ఉపయోగించడం కొనసాగించారు.

\end{document}
